% theorems.tex — derived theorems and the Y-fixed-point.

\section{Theorems}

\subsection{Y-fixed-point: math-coding applies to itself}

\begin{theorem}[y-fixed-point.math-coding]
math-coding is a Y-fixed point:
$$Y(\mathrm{math\text{-}coding}) = \mathrm{math\text{-}coding}$$
The convention, applied to itself, yields itself.
\end{theorem}

\begin{proof}
math-coding is a function $\mathcal{M} : \mathrm{Project} \to
\mathcal{P}(\Dec)$. Applied to itself, it produces:
\begin{enumerate}
\item the five foundations as ordinary packets
      (\texttt{curry-howard}, \texttt{temporal}, \texttt{constructive},
      \texttt{categorical}, \texttt{motivation}),
\item the three extensions as ordinary packets
      (\texttt{process-fsm}, \texttt{dialectic-tas},
      \texttt{actor-discipline}),
\item the kernel $S$ that type-checks those packets using the same
      code that checks user packets.
\end{enumerate}
All three are themselves packets verified by the same kernel. There
is no special path for the foundations; this is the structural
expression of the Y-fixed-point property. If any proposition in a
foundation packet is false, the kernel detects drift in that
foundation just as it does for a user packet.
\end{proof}

\subsection{Drift detection}

\begin{theorem}[drift.detection]
$$L(d, r) = \mathrm{Drift} \iff d.\mathtt{witness} = \text{sha} \;\wedge\; \mathrm{prop}(r, \text{sha}, d.\mathtt{name}) \neq d.\mathtt{proposition}$$
That is, drift occurs exactly when the proposition at the witness
commit differs from the current proposition.
\end{theorem}

\begin{proof}
If $\mathrm{prop}(r, \text{sha}, d.\mathtt{name}) = d.\mathtt{proposition}$,
the proposition is consistent with history and $L = \mathrm{Applied}$
(possibly $\mathrm{Stale}$ if files changed). If they differ, the
proposition has been modified without a supersession, which is the
definition of drift. The kernel implements this comparison through
\texttt{git show sha:math/name/packet.md}.
\end{proof}

\subsection{SPO transitivity preserves lineage}

\begin{theorem}[spo-transitivity]
If $d_1 \preceq d_2$ and $d_2 \preceq d_3$, then $d_1 \preceq d_3$
through the composition of two supersession chains.
\end{theorem}

\begin{proof}
By definition of transitivity in
\texttt{categorical.supersession-spo}, the composition of two
supersession links is a supersession link. The kernel preserves
chains through the \texttt{mathc graph} command, which renders the
Hasse diagram of the partial order.
\end{proof}

\subsection{Kernel uniformity}

\begin{theorem}[kernel-uniformity]
The kernel $S$ applies uniformly to all decisions, including the
foundations and extensions. No special path exists for axiomatic
packets.
\end{theorem}

\begin{proof}
\texttt{core/check.ml} iterates over \texttt{math/} recursively and
applies \texttt{Check.check : decision -> verdict} to every packet
without distinction. The foundations satisfy the schema (V1--V7)
the same way user packets do. This is the operational expression of
the Y-fixed-point: the kernel that verifies packets also verifies
the description of the kernel.
\end{proof}